\chapter{Podatki in predobdelava}
\label{sec:podatki}

\section{Podatkovna zbirka \MAHNOB{}}
\label{subsec:podatkovna-zbirka}

V delu uporabljamo podatke iz podatkovne zbirke \MAHNOB{}~\cite{soleymani2012mahnob}.
Gre za večmodalno zbirko, namenjeno raziskavam prepoznavanja čustev in implicitnega označevanja multimedijskih vsebin.
Zbirka vsebuje dva glavna eksperimentalna sklopa: prepoznavanje čustvenih stanj iz odzivov na čustvene videoposnetke in implicitno označevanje na podlagi odzivov na pravilne oziroma napačne oznake pod sliko ali posnetkom.
V tej nalogi uporabljamo samo prvi sklop, to je eksperiment vzbujanja čustev z videoposnetki.

Izvirni članek poroča o 30 rekrutiranih udeležencih, od katerih je bilo v analizi uporabljenih 27.
Udeleženci P9, P12 in P15 so bili izključeni zaradi tehničnih težav oziroma nepopolnega zajema podatkov.
V tej nalogi po vzoru izvirnega članka izključimo iste tri udeležence.
V eksperimentu vzbujanja čustev so udeleženci gledali 20 čustveno vzbujajočih videoposnetkov.
Pred vsakim čustvenim posnetkom je bil prikazan še kratek nevtralni posnetek, odziv nanj pa je bil shranjen ločeno.
Po ogledu posameznega posnetka so udeleženci podali samoocene čustvene oznake, vzburjenosti, valence, dominantnosti in predvidljivosti.

Čeprav je \MAHNOB{} večmodalna zbirka, v tej nalogi uporabljamo izključno signale sledilnika pogleda.
Uporabljeni signali so podrobneje opisani v razdelku~\ref{subsec:vhodni-signali}.
Meritve so bile zajete s sledilnikom Tobii X120 pri frekvenci 60~Hz.

\section{Predhodni poskus z zbirko \ESEED{}}
\label{subsec:eseed}

Pred prehodom na \MAHNOB{} je bil del razvoja izveden na podatkovni zbirki \ESEED{}.
Ta del ni osrednji eksperimentalni prispevek naloge, temveč služi kot zgodovinska motivacija za spremembo podatkovne zbirke.
Začetni poskusi so pokazali težave z zanesljivostjo oznak, kakovostjo signalov in stabilnostjo učenja.
Zato smo glavno eksperimentalno evalvacijo izvedli na podatkovni zbirki \MAHNOB{}, ki ima za naš primer uporabnejšo množico signalov sledilnika pogleda.

\section{Ciljne oznake}
\label{subsec:ciljne-oznake}

Oznake v tej nalogi izhajajo iz udeleženčevih samoocen čustvenega razreda po ogledu videoposnetka.
Diskretne čustvene razrede preslikamo v razrede valence po preslikavi iz izvirnega članka~\cite{soleymani2012mahnob}, ki je prikazana v dodatku~\ref{app:preslikava-valence}.
S tem najprej dobimo tri razrede: neprijetno, nevtralno in prijetno valenco.
Nevtralni razred odstranimo, da ciljna naloga vsebuje izrazitejši kontrast med razredoma in s tem modelom omogočimo boljše pogoje za učenje vzorcev v podatkih.

Oznaka valence je določena na ravni kombinacije udeleženca in videoposnetka.
Vsa časovna okna in s tem vsi grafi istega ogleda videoposnetka zato prejmejo isto oznako.

Vzburjenosti v tej nalogi ne klasificiramo.
Razlog je predvsem večja šumnost oznak v uporabljeni podatkovni zbirki.
Kot kaže matrika zmede v dodatku~\ref{app:ujemanje-oznak-s-samoocenami}, se samoocene stopnje vzburjenosti slabše ujemajo z definiranimi razredi kot samoocene valence.


\begin{figure}[t]
    \centering
    \includegraphics[width=0.82\textwidth]{porazdelitev_razredov_stevila_retain_2026-06-12.png}
    \caption{Porazdelitev 10-sekundnih oken po razredih binarne valence.}
    \label{fig:porazdelitev-valenca-stevila}
\end{figure}

\END_NOTE{Slika je delovna verzija iz runa RETAIN_2026-06-12_16-29-08. Pred oddajo preveri, ali je generirana iz končnega protokola, in jo po potrebi regeneriraj v paleti B.}

\section{Priprava vhodnih signalov}
\label{subsec:vhodni-signali}

\subsection{Skupine signalov}

Tukaj navedemo konkretne skupine signalov, ki jih uporabimo v eksperimentih.
V uporabljeni podatkovni zbirki imamo na voljo pet osnovnih vrst signalov. To so:
\begin{itemize}
    \item časovni žig meritve $t_i$,
    \item koordinati pogleda na zaslonu $(x_i,~y_i)$ -- povprečje med meritvama levega in desnega očesa,
    \item velikost leve in desne zenice $(p_i^{(l)},~p_i^{(r)})$,
    \item oddaljenost oči od sledilnika pogleda $d_i$ ter
    \item informacija o fiksacijah, tj. pripadnost fiksaciji ter trajanje fiksacije.
\end{itemize}

Večina podatkovnih zbirk s podatki sledilnika pogleda vsebuje bodisi le koordinate pogleda bodisi le velikosti zenic bodisi oboje.
Nekatere zbirke pa ponujajo tudi dodatne signale, kot sta v našem primeru oddaljenost od sledilnika ter informacije o fiksacijah.
Da bi s primerjalno študijo pokrili čim širši razpon različnih tipov zbirk, \MAHNOB{} podatke razdelimo v štiri delno prekrivajoče se eksperimentalne množice signalov.
V vseh štirih množicah uporabimo tudi informacijo o času.
Množica \emph{samo pogled} poleg časa vsebuje koordinati pogleda, množica \emph{samo zenici} vsebuje velikosti leve in desne zenice, množica \emph{pogled in zenici} združi koordinati pogleda in velikosti zenic, množica \emph{vsi signali} pa doda še oddaljenost od sledilnika pogleda in informacijo o fiksacijah.
Podrobna povezava med temi množicami, značilkami vozlišč in razpoložljivimi tipi povezav je podana pri opisu grafovske konstrukcije v razdelku~\ref{subsec:grafovska-konstrukcija-po-mnozicah}.
\END_NOTE{Z zunanjo pomočjo mentorja ali slovenista najdi/potrdi izraze in jih potem uporabljaj konsistentno v diplomi. Posebej preveri, ali naj bodo končna imena množic \emph{samo pogled}, \emph{samo zenici}, \emph{pogled in zenici} in \emph{vsi signali}, ter nato ta imena uporabi enako v besedilu, tabelah in slikah.}


\subsection{Čiščenje podatkov}
\label{subsec:ciscenje}

Podatke pred uporabo očistimo.
Najprej iz lokalne kopije zbirke \MAHNOB odstranimo podatke, ki jih v tej nalogi ne uporabljamo, in ohranimo le podatke eksperimenta vzbujanja čustev.
Nato odstranimo vrstice brez signalov pogleda ali velikosti zenic.
Na voljo imamo tudi signale veljavnosti, na podlagi katerih odstranimo meritve, za katere je sledilnik poročal, da je meritev za vsaj eno oko neveljavna.
Dodatno uporabimo filter osamelcev \texttt{q01}-\texttt{q99}, s katerim odstranimo izrazite ekstreme, ki izkrivljajo porazdelitve signalov.
Na koncu kratke vrzeli v numeričnih signalih interpoliramo, signala velikosti leve in desne zenice pa dodatno zgladimo s kratkim drsečim povprečjem.

\section{Normalizacija}
\label{subsec:normalizacija}

Posamezne skupine signalov so izražene v različnih velikostnih redih.
Modeli strojnega učenja, kot so \ModelName{SVM}, \ModelName{MLP} ter GNN, so občutljivi na merilo značilk.
Zato signale standardiziramo.
Povprečje in standardni odklon izračunamo na učni podmnožici, nato pa na validacijski in testni podmnožici standardizacijo uporabimo z istimi parametri kot na učni.
Pri grafovskih modelih na enak način standardiziramo tudi značilke povezav, kadar jih uporabljena množica signalov vsebuje.

\section{Segmentacija v časovna okna}
\label{subsec:windowing}

Časovno vrsto razdelimo na neprekrivajoča se časovna okna dolžine $T=10~s$.
Vsako uporabno časovno okno predstavlja en učni primer.
Za $10~s$ smo se odločili na podlagi sorodnega dela na podatkih sledilnika pogleda, kjer so se daljša okna med dolžinami $1$, $3$, $5$ in $10~s$ praviloma izkazala za bolj informativna~\cite{bozak2024feature}.
Daljših oken ne uporabimo, ker bi težje predpostavili, da celotno okno opisuje enoten čustveni odziv, saj se lahko ta med ogledom posnetka že opazno spremeni~\cite{zhang2020rcea}.
Tudi sorodna dela na prepoznavanju čustev iz očesnih signalov pogosto uporabljajo krajša okna, na primer $1$-$2~s$~\cite{wang2018automatic,wu2021emotion}.

Oznaka okna je določena s parom udeleženec-posnetek, zato so vsa okna istega ogleda videoposnetka označena z isto čustveno oznako.
Uporabnost okna določimo glede na kakovost podatkov v oknu.

Pri vzorčni frekvenci 60~Hz bi 10-sekundno okno brez manjkajočih vrednosti vsebovalo 600 vzorcev.
Zaradi čiščenja, odstranjevanja neveljavnih vzorcev in filtriranja je dejansko število vzorcev v oknih manjše.
Povprečno število meritev v oknu po čiščenju znaša \TODO{vstavi številko}.

Po tvorbi časovnih oken se cevovod razdeli glede na uporabljeno predstavitev.
Negrafovski osnovni modeli iz okna dobijo agregirane tabelarične značilke, zamrznjeni prednaučeni predstavitveni modeli neposredno prejmejo časovne vrste izbranih signalov sledilnika pogleda, grafovski modeli pa grafovsko predstavitev okna.
Podrobnosti teh predstavitev so opisane v poglavjih~\ref{sec:grafovska-predstavitev} in~\ref{sec:modeli}.
V tem poglavju opišemo le skupne korake priprave podatkov.

\section{Povzetek predobdelave}

Tabela~\ref{tab:mahnob-predobdelava} povzema, kako se skozi opisane korake predobdelave zmanjša količina podatkov, uporabljena v končni eksperimentalni primerjavi.

\begin{table}[t]
    \centering
    \caption{Zmanjševanje količine podatkov od izvirne zbirke \MAHNOB{} do končnih časovnih oken, uporabljenih v eksperimentalnem cevovodu.}
    \label{tab:mahnob-predobdelava}
    \small
    \begin{tabularx}{\textwidth}{@{}>{\raggedright\arraybackslash}p{0.28\textwidth} c c >{\raggedright\arraybackslash}X@{}}
        \toprule
        Korak & Vrstice & Udeleženci & Opis \\
        \midrule
        Izbira izvirne zbirke \MAHNOB{} & -- & 30 & Večmodalna zbirka; izvirni članek v analizi uporabi 27 udeležencev. \\
        Izbira podmnožice s čustvenimi oznakami & 3\,797\,165 & 24 & Ohranimo le posnetke iz eksperimenta vzbujanja čustev, ki imajo podatke sledilnika pogleda in čustvene oznake. \\
        Izključitev "slabih" udeležencev & 3\,535\,842 & 22 & Izključeni P9, P12 in P15. \\
        Izbira veljavno označenih vrstic & 2\,995\,632 & 22 & Ohranjeni so samo označeni podatki, uporabljeni za učenje modelov čustvenih stanj. \\
        Priprava podatkov za učenje & 2\,814\,005 & 22 & Odstranjene so vrstice brez signalov pogleda ali zenic. \\
        Filter osamelcev \texttt{q01}-\texttt{q99} & 2\,639\,048 & 22 & Robustno filtriranje štirih osnovnih signalov pogleda in zenic. \\
        Tvorba uporabnih 10-sekundnih časovnih oken & 5\,158 & 22 & Vsako okno postane en učni primer oziroma graf. \\
        \bottomrule
    \end{tabularx}
\end{table}
